% !TEX root = ./ECEN-601 - Homework 3 - Brody Jordan.tex

\documentclass[11pt]{article}

\usepackage[
    assignmentDueDate={September 26, 2025 at 11:59 pm}
]{C:/Users/bljor/Sync/Courses/assignment}
\newcommand\op{\mathop{\circ}\nolimits}

\begin{document}
\makeAssignmentTitle

% —————— PROBLEM 1 —————— %
\begin{homeworkProblem}
    Prove that if a Cauchy sequence in a metric space (not assumed to be complete) has a limit
    point, then it has a limit. [A point $x$ is a limit point of $\{x_n\}$ if every neighborhood
            of $x$ contains infinitely many elements of the sequence.] \\

    \solution

\end{homeworkProblem}

\pagebreak

% —————— PROBLEM 2 —————— %
\begin{homeworkProblem}
    Call two metrics $d_1$ and $d_2$ on the same set $M$ equivalent if there exist positive constants
    $c_1$, $c_2$ such that $d_1(x, y) \leq c_2 d_2(x, y)$ and $d_2(x, y) \leq c_1 d_1(x, y)$ for all $x$ and $y$ in $M$. Prove
    that $x_n \to x$ in the $d_1$-metric if and only if $x_n \to x$ in the $d_2$-metric, if $d_1$ and $d_2$ are equivalent. \\

    \solution
\end{homeworkProblem}

\pagebreak

% —————— PROBLEM 3 —————— %
\begin{homeworkProblem}
    If $f : M \rightarrow N$ and $g : N \rightarrow P$ are continuous, prove $g \hspace{0.05em} \circ \hspace{0.05em} f : M \rightarrow P$ is continuous. Given an
    example where $g \hspace{0.05em} \circ \hspace{0.05em} f$ is continuous but $g$ and $f$ are not. \\

    \solution
\end{homeworkProblem}

\pagebreak

% —————— PROBLEM 4 —————— %
\begin{homeworkProblem}
    Let $a$ be a fixed real number with $1 < a < 3$. Prove that the mapping $f (x) = (x/2) + (a/2x)$
    satisfies the hypotheses of the contractive mapping principle on the domain $(1, \infty)$. What is
    the fixed point? \\

    \solution
\end{homeworkProblem}

\pagebreak

% —————— PROBLEM 5 —————— %
\begin{homeworkProblem}
    If $f$ satisfies the hypotheses of the contractive mapping principle and $x_1$ is any point in $M$,
    show that $d(x_1, x_0) \leq d(x_1, f (x_1))/(1 - r)$ where $x_0$ is the fixed point. Informally, this says
    that if $f (x_1)$ is close to $x_1$, then $x_1$ is close to the fixed point (but $r$ must not be too close to
    1 for this to be a good estimate). \\

    \solution
\end{homeworkProblem}
\end{document}